\section{Delivering Frame-Accurate Results in the Browser}
\label{sec:player}

Tracking results are indexed by frame, so detections, transforms, pose anchors, and
annotations must refer to the image actually displayed. An ordinary HTML video
element exposes media time rather than stable frame identity, which is insufficient
for frame stepping, hit-testing, and inspection across an occlusion.

The WebCodecs path~\cite{webcodecs} discards packets with negative presentation
timestamps (PTS), then orders the remainder by PTS and packet sequence number. Position
in that list is the canonical zero-based frame index. Preview seeks from a preceding
keyframe and decodes forward, reusing a bounded cache for nearby frames. A result is
published only if it still belongs to the latest seek.

Export decodes independently, advances through the same packet-index convention, and
resynchronizes to the closest PTS if necessary. The recovered index selects the rider
track. Overview preview additionally reads the global stabilization transform, whereas
focused preview and export share the pose-anchor crop function. Decoded images and
output surfaces remain separate. This machinery is not a tracking contribution; it
makes the result inspectable and exportable without silently shifting an overlay or
crop to a neighboring frame.

\begin{figure}[t]
  \centering
  \includegraphics[width=\linewidth]{player-frame-accurate-playback.pdf}
  \caption{Preview and export decode independently but identify frame $i$ by the same
  ordered-packet index. Preview seeks from a preceding keyframe; export advances its
  own decoder and can resynchronize by timestamp. Both use the index for track lookup,
  while only focused rendering shares the pose-anchor crop function.}
  \label{fig:player-frame-contract}
\end{figure}
